\documentclass[10pt,a4paper]{amsart}
\usepackage[margin=19mm]{geometry}
\usepackage{amsmath,amssymb,mathtools}
\usepackage[scr=rsfs,cal=euler]{mathalfa}
\usepackage{xfrac}
\usepackage[unicode,hidelinks,bookmarks=false]{hyperref}
\newcommand{\ie}{i.e.\ }
\newcommand{\cf}{cf.\ }
\newcommand{\resp}{respectively\ }
\newcommand{\Subsection}{\subsection}
\providecommand{\nobreakdash}{\nobreak\hbox{-}}
\setlength{\emergencystretch}{2em}
\multlinegap0pt
\usepackage{amssymb}
\usepackage{xcolor}
\newtheorem{lemma}{Lemma}
\newtheorem{remark}{Remark}
\newtheorem{proposition}{Proposition}
\newcommand{\F}{\mathbb{F}}
\newcommand{\Tr}{{\rm Tr}}
\newcommand{\V}{\mathbb{V}}
\renewcommand{\to}{\rightarrow}
\title{On Generalizations of Maiorana--McFarland and $\mathcal{PS}_{ap}$ Functions}
\author{Sezel Alkan, Nurdag\"ul Anbar, Athina Avrantini, Erroxe Etxabarri-Alberdi, Tekg\"ul Kalayc{\i}, Beatrice Toesca}
\date{Source: arXiv:2603.28485v1 (CC BY 4.0)}
\begin{document}
\maketitle

\noindent\emph{Source and licence.} On Generalizations of Maiorana--McFarland and $\mathcal{PS}_{ap}$ Functions. Version arXiv:2603.28485v1 (CC BY 4.0), distributed by arXiv under the Creative Commons Attribution 4.0 International licence, http://creativecommons.org/licenses/by/4.0/, as recorded in the arXiv licence metadata for this exact version and shown on the abstract page https://arxiv.org/abs/2603.28485v1. Full licence notice: \texttt{licenses/arxiv-2603.28485-CC-BY-4.0.txt}.\par\smallskip

\noindent\textit{Editorial scope.} Counted unit: the article's proposition thm:main (source0.tex lines 711--733). The article's complete original proof of it is delivered with it (source0.tex lines 736--908, one \texttt{proof} environment of 173 lines). No mathematics is written, repaired or re-derived: the statement and the proof are the article's own lines, in the article's own notation, and no retained line is substituted or re-flowed.  Retained uncounted as the source-local context the proof needs: lines 585--598; lines 632--637; lines 651--663.  Nothing was omitted from any retained span, so the package carries no omission record.  No retained line contains a citation, so the wrapper carries no bibliography and no bibliography entry.  No retained line contains a figure environment, an image-inclusion command or drawing code, so no image is delivered and the package declares no asset dependency.  Editorial formatting only, in addition: an AMS \texttt{amsart} preamble that restates the article's own declarations and macros used by the retained lines, namely the environments \texttt{lemma}, \texttt{remark}, \texttt{proposition} on separate counters, together with the standard packages those lines need.  The retained environments print in this document as Lemma 1, Remark 1, Proposition 1 in order of appearance, whereas the article prints the same items with the numbering of its own full document.  The complete native source of the article as distributed in the arXiv e-print for this exact version is bundled unchanged as \texttt{sources/197447/source0.tex}.
	\begin{lemma}\label{lem:second}
		Let \(f \colon \V_n^{(2)} \to \F_2\) be a Boolean function, and let
		\(a,b \in \V_n^{(2)}\) be linearly independent. Let
		\(L \colon \V_n^{(2)} \to \V_n^{(2)}\) be a linear permutation, and let
		\(c,d \in \V_n^{(2)}\) and \(e \in \F_2\).
		Then \(D_a D_b f(x)\) is identically equal to \(0\) (respectively, identically equal to $1$)
		for all \(x \in \V_n^{(2)}\) if and only if
		\begin{equation}\label{eq:second}
			D_{L^{-1}(a)} D_{L^{-1}(b)} \left(\ f\bigl(L(x)+c\bigr)
			+ \langle d,x\rangle_n + e \ \right) 
		\end{equation}
		is also identically \(0\) (respectively, identically \(1\)) for all
		\(x \in \V_n^{(2)}\).
	\end{lemma}

	\begin{remark}\label{rem:fcircL}
		Let \(f \colon \V_n^{(2)} \to \F_2\) be a Boolean function, and let
		\(L \colon \V_n^{(2)} \to \V_n^{(2)}\) be a linear permutation. Then, by Lemma~\ref{lem:second},
		\(\mathcal{W} \subseteq \V_n^{(2)}\) is an \(\mathcal{M}\)-subspace of \(f\) if and only if
		\(L^{-1}(\mathcal{W})\) is an \(\mathcal{M}\)-subspace of \(f \circ L\).
	\end{remark}

	\begin{align}\label{eq:secondder}
		&D_{\nu_1}D_{\nu_2}f(x,y,z)
		= f(x,y,z)
		+ f(x+u_1,y+v_1,z+w_1)
		+ f(x+u_2,y+v_2,z+w_2)\\ \nonumber
		&\qquad\qquad 
		+ f(x+u_1+u_2,y+v_1+v_2,z+w_1+w_2)\\[0.5em] \nonumber
		&=f^{(z)}(x)
		+ f^{(z+w_1)}(x+u_1)
		+f^{(z+w_2)}(x+u_2)
		+ f^{(z+w_1+w_2)}(x+u_1+u_2)
		+ \Tr_1^k\!\bigl(w_1v_2+w_2v_1\bigr).
	\end{align}

	\begin{proposition}\label{thm:main}
		Let \(h\) be an \(MM\) function on \(\V_n^{(2)}\) 
		having a unique 
		\(\mathcal{M}\)-subspace \(\mathcal{W}\) of dimension \(n/2\), such that if $\mathcal{U}$ is a nontrivial
		\(\mathcal{M}\)-subspace of $h$, then $\mathcal{U}\subseteq \mathcal{W}$. Let  \(\widetilde{\mathcal{W}}\) be a complementary subspace of \(\V_n^{(2)}\) satisfying 
		\(\V_n^{(2)} =\mathcal{W} \oplus \widetilde{\mathcal{W}}\) (i.e., $\mathcal{W} \cap \widetilde{\mathcal{W}}=\{0\}$ and \(\V_n^{(2)} =\mathcal{W} + \widetilde{\mathcal{W}}\)), and let \(L\) be a linear permutation of \(\V_n^{(2)}\) with 
		\({L}( \widetilde{\mathcal{W}} ) = \mathcal{W}\).
		
		For each \(z\in\mathbb{F}_{2^k}\), define  
		\begin{align}\label{eq:g}
			f^{(z)}(x) =
			\begin{cases}
				h(x), & \text{if } z \in V,\\[6pt]
				(h\circ L)(x), & \text{if } z \not\in V,
			\end{cases}
		\end{align}
		where $	V = \bigl\{\, z \in \mathbb{F}_{2^k} : \Tr_1^k(z)=0 \,\bigr\}$.
		
		With this choice of \(f^{(z)}\) in \eqref{eq:g}, if \(n > 2k+4\), 
		then the function $	f:\V_n^{(2)} \times \mathbb{F}_{2^k} \times \mathbb{F}_{2^k} \to \mathbb{F}_{2}$, defined by $f(x,y,z) = f^{(z)}(x) + \Tr_1^k\!\bigl(y z\bigr)$, is not (equivalent to) an \(MM\) function. 
		Equivalently, if \(\mathcal{U}\) is an \(\mathcal{M}\)-subspace of \(f\), then 
		\(\dim(\mathcal{U}) < n/2 + k\).
	\end{proposition}

	\begin{proof}
		We recall that, by Remark~\ref{rem:fcircL}, the subspace
		\(L^{-1}(\mathcal{W}) = \widetilde{\mathcal{W}}\) is an \(\mathcal{M}\)-subspace of
		\(h \circ L\) of dimension \(n/2\).
		Moreover, by assumption, \(\widetilde{W}\) is the unique
		\(\mathcal{M}\)-subspace with the property that every nontrivial
		\(\mathcal{M}\)-subspace of \(h\circ L\) is contained in \(\widetilde{W}\).
		Equivalently, for {distinct} nonzero
		\({a},{b}\in \V_n^{(2)}\), the equality $D_{{a}}D_{{b}} \, h (L (x))=0$ holds for all $x\in \V_n^{(2)}$
		if and only if \({a},{b}\in \widetilde{W}\).
		
		\medskip
		
		Let \(\mathcal{U}\) be an \(\mathcal{M}\)-subspace of \(f\).
		We derive an upper bound on the dimension \(\dim(\mathcal{U})\) of $\mathcal{U}$ in two steps.
		
		\medskip
		\noindent\textbf{Step 1.}
		Define
		\begin{equation}\label{eq:UV}
			\mathcal{U}_V
			=
			\bigl\{\,
			(u,v,w)\in \mathcal{U} : \Tr_1^k(w)=0
			\,\bigr\}.
		\end{equation}
		Let
		\(\nu_1=(u_1,v_1,w_1)\) and \(\nu_2=(u_2,v_2,w_2)\)
		be distinct nonzero elements of \(\mathcal{U}_V\).
		Since \(V\) is a subspace, \(w_1+w_2\in V\).
		
		Assume first that
		\(\Tr_1^k(w_1v_2+w_2v_1)=0\).
		Then for any \(z\in V\),
		Equation~\eqref{eq:secondder} yields
		\begin{align*}
			D_{\nu_1}D_{\nu_2}f(x,y,z)
			&= h(x)+h(x+u_1)+h(x+u_2)+h(x+u_1+u_2) \\
			&= D_{u_1}D_{u_2}h(x).
		\end{align*}
		Hence,
		\(D_{\nu_1}D_{\nu_2}f(x,y,z)=0\) for all \(x\in \V_n^{(2)}\)
		if and only if either
		\(u_1,u_2\) are distinct nonzero elements of the
		\(\mathcal{M}\)-subspace of \(h\), or $(u_1,u_2)\in \{(0,u),(u,0),(u,u):u\in \V_n^{(2)}\}$.
		
		Similarly, for \(z\notin V\), $D_{\nu_1}D_{\nu_2}f(x,y,z)
		= D_{u_1}D_{u_2}\, h\circ L(x)$,
		and vanishing occurs if and only if either
		\(u_1,u_2\) belong to the \(\mathcal{M}\)-subspace of
		\(h\circ L\), or  $(u_1,u_2)\in \{(0,u),(u,0),(u,u):u\in \V_n^{(2)}\}$.
		Since \(\mathcal{W}\cap \widetilde{\mathcal{W}}=\{0\}\),
		we must have $(u_1,u_2)\in \{(0,u),(u,0),(u,u):u\in \V_n^{(2)}\}$.
		
		We define $\mathfrak S \subseteq \mathcal U_V$ to be a subset such that, 
		for any $(u_i,v_i,w_i), (u_j,v_j,w_j)\in \mathfrak S$, one has $\Tr_1^k\!\left(w_i v_j + w_j v_i\right)=0$. In other words,
		\[
		\mathfrak{S}
		=
		\left\{
		\nu=(u,v,w)\in \mathcal{U}_V
		\; :\;
		\Tr_1^k\!\left(w v_i + w_i v\right)=0
		\ \text{for all } (u_i,v_i,w_i)\in \mathfrak{S}
		\right\}.\]
		Choose \(\mathfrak{S}\) of maximal dimension.
		If \(\mathfrak{S}\) contains two elements with nonzero first components,
		then these components must be the same.
		Hence, for a fixed nonzero \(u\in \V_n^{(2)} \),
		\(\mathfrak{S}\) is contained in
		\[
		\{0,u\}\times
		\left\{
		(v_j,w_j)\in \mathbb{F}_{2^k}^2
		\; : \;
		\Tr_1^k(w_j)=0
		\ \text{and}\
		\Tr_1^k(w_j v_i + w_i v_j)=0
		\ \text{for all } i
		\right\},
		\]  
		which is an \(\mathcal{M}\)-subspace of \(f\). 
		
		\smallskip
		
		\noindent\textbf{Case (i):}
		If \(\mathfrak{S}=\{(0,0,0)\}\), then
		we claim that \(\dim(\mathcal{U}_V)\leq 2\).
		Suppose, for the sake of contradiction, that \(\mathcal{U}_V\) contains
		three linearly independent vectors
		\(\nu_i=(u_i,v_i,w_i)\), \(i=1,2,3\).
		Since \(\mathfrak{S}=\{(0,0,0)\}\), we have $\Tr_1^k\!\bigl(w_i v_j + w_j v_i\bigr)=1$ for all $i\neq j$.
		By linear independence, the vectors
		\((u_1+u_2,v_1+v_2,w_1+w_2)\) and
		\((u_3,v_3,w_3)\)
		are distinct nonzero elements of \(\mathcal{U}_V\).
		Moreover, we have $\Tr_1^k\!\bigl((w_1+w_2)v_3+w_3(v_1+v_2)\bigr)=0$.
		Hence,
		\((u_1+u_2,v_1+v_2,w_1+w_2)\) and
		\((u_3,v_3,w_3)\)
		both belong to \(\mathfrak{S}\),
		contradicting the assumption that
		\(\mathfrak{S}=\{(0,0,0)\}\).
		
		\smallskip
		\noindent\textbf{Case (ii):}
		Suppose that \(\mathfrak{S}\neq\{(0,0,0)\}\). If
		\(\mathfrak{S} =\mathcal{U}_V\),
		then $\mathrm{dim}(\mathfrak{S}) =  \mathrm{dim}(\mathcal{U}_V)\leq 2k$.
		
		In the case $\mathfrak{S} \not\subseteq \mathcal{U}_V$, choose
		\(\nu_1=(u_1,v_1,w_1)\in \mathcal{U}_V\setminus \mathfrak{S}\).
		By maximality, there exists
		\(\nu_2=(u_2,v_2,w_2)\in \mathfrak{S}\)  with
		\(\Tr_1^k(w_1v_2+w_2v_1)=1\),
		which forces \(u_2=u\) and \(u_1\notin\{0,u\}\).
		Otherwise,
		\[
		h(x)+h(x+u_1)+h(x+u_2)+h\bigl(x+u_1+u_2\bigr)=0,
		\]
		which would imply that, for any \(z\in V\), $D_{\nu_1}D_{\nu_2}f(x,y,z)=1$.
		
		Conversely, for any \(\nu_i=(u_i,v_i,w_i)\in \mathfrak{S}\) with \(u_i=u\),
		we have $\Tr_1^k\!\bigl(w_1v_i+w_iv_1\bigr)=1$.
		Indeed, if this were not the case, then
		\(\{u,u_1\}\subseteq \mathcal{W} \cap \widetilde{\mathcal{W} }\),
		which contradicts our assumption that
		\(\mathcal{W} \cap \widetilde{\mathcal{W} }=\{0\}\).
		
		We claim that any $\mathcal{M}$-subspace of $f$ contains at most two vectors whose last components have trace zero and whose first components are linearly independent.
		This implies that
		\begin{align*}
			\mathcal{U}_V
			\subseteq
			\langle u_1,u_2\rangle
			\times
			\Bigl\{(v,w)\in \mathbb{F}_{2^k}\times \mathbb{F}_{2^k} :
			\Tr_1^k(w)=0
			\Bigr\},
		\end{align*}
		which is a subspace of dimension \(2k+1\).
		Suppose, for the sake of contradiction, that there exist three vectors \(\nu_i=(u_i,v_i,w_i) \in \mathcal{U}_V\), \(i=1,2,3,\) with linearly independent first components satisfying $\Tr_1^k\!\bigl(w_i v_j+w_j v_i\bigr)=1$ for all $i\neq j$.
		Since \(\{u_1,u_2,u_3\}\) is linearly independent, the vectors
		\((u_1+u_2,v_1+v_2,w_1+w_2)\) and
		\((u_3,v_3,w_3)\)
		are distinct nonzero elements of $\mathcal{U}_V$. 
		Moreover, we have $\Tr_1^k\!\bigl((w_1+w_2)v_3+w_3(v_1+v_2)\bigr)=0$.
		By the argument above, this implies that
		\[
		(u_1+u_2,u_3)\in \{(0,u),(u,0),(u,u):u\in \V_n^{(2)}\}.
		\]
		Since \(u_1\neq u_2\) and \(u_3\neq 0\), we must have \(u_1+u_2=u_3\),
		which contradicts the linear independence of \(\{u_1,u_2,u_3\}\).
		
		Thus, any \(\mathcal{M}\)-subspace of \(f\) consisting of 
		\((u,v,w)\) with \(w\in V\) has dimension at most \(2k+1\). In particular, the dimension of $\mathcal{U}_V$ is at most \(2k+1\).
		
		\medskip
		\noindent\textbf{Step 2:}  
		We now consider the set $\bigl\{\, (u,v,w) \in \mathcal{U} : \Tr_1^k(w)=1 \,\bigr\}$. 
		Equivalently, for a fixed $(u_1,v_1,w_1)\in\mathcal{U} $ with $\Tr_1^k(w_1)=1$, we can consider  
		\[
		\bigl\{\, (u_1,v_1,w_1),\; (u_1+u,\;v_1+v,\;w_1+w) \, : \, (u,v,w)\in \mathcal{U}, \, \Tr_1^k(w)=1 \,\bigr\},
		\]
		as both sets span the same subspace.
		Since \(\Tr_1^k(w_1)=\Tr_1^k(w)=1\), we have \(\Tr_1^k(w_1+w)=0\).  
		In particular,  
		\((u_1+u,v_1+v,w_1+w)\in \mathcal{U}_V\), 
		where \(\mathcal{U}_V\) is defined in Equation \eqref{eq:UV}.  
		Hence, \(\mathcal{U}_V\) and \((u_1,v_1,w_1)\) generate the whole space \(\mathcal{U}\), which implies that 
		\(\dim (\mathcal{U})\leq 2k+2\). Then our assumption $n > 2k + 4$ implies that $\dim(\mathcal U) < \tfrac{n}{2} + k$,
		which yields the desired conclusion.
	\end{proof}

\end{document}
